\begin{exercise}[The general spin-statistics connection]{I-CH02-006}
Let a nonzero local field multiplet $\Phi_a(x)$ in $D=4$ transform in the finite
irreducible Lorentz representation $(A,B)$. Assume Poincar\'e covariance, a
unique invariant vacuum, a positive Hilbert-space metric, the spectrum
condition, and cyclicity of the vacuum for the local field algebra. Its
restriction to rotations contains spins
\[
  j=|A-B|,|A-B|+1,\ldots,A+B.
\]
First show that every spin in this list has
$(-1)^{2j}=(-1)^{2(A+B)}$.

For space-like separated arguments write graded locality as
\[
  \Phi_a(x)\Phi_b{}^\dagger(y)
  =\eta\,\Phi_b{}^\dagger(y)\Phi_a(x),
  \qquad \eta\in\{+1,-1\}.
\]
Use the Wightman spin-statistics lemma: the spectrum condition analytically
continues vacuum correlation functions to the forward and backward tubes; at
a space-like Jost configuration, continuation through the connected complex
Lorentz group reverses the two fields and contributes the central-rotation
factor $(-1)^{2(A+B)}$. Explain this continuation and show that a nonzero field
requires
\[
  \eta=(-1)^{2j}.
\]
State where positivity and locality enter, check scalar, spinor, and vector
fields, and extend the conclusion to space-time dimension $D>4$ by using the
central $2\pi$ rotation in $\operatorname{Spin}(D-1)$. Explain why braid
statistics in lower dimension evades this argument.
\end{exercise}
